\documentclass[12pt]{article}
\usepackage{graphicx}
\usepackage[margin=1in]{geometry}
\usepackage{caption}
\usepackage{xcolor}
\usepackage{float}
\usepackage{amsmath}
\usepackage{array}
\usepackage{multicol}
\usepackage{hyperref}
\usepackage{url}

\setlength{\columnsep}{1cm}

\begin{document}

% Title Block
\begin{figure}[H]
    \begin{minipage}{0.45\textwidth}
        \includegraphics[width=\textwidth]{iiitb_comet.jpeg} % <-- Replace with actual image
    \end{minipage} \hfill
    \begin{minipage}{0.45\textwidth}
        \textbf{Name: S.Rithrishareddy} \\
        \textbf{Batch: COMETFWC028} \\
        \textbf{Date: 06 July 2025}
    \end{minipage}
\end{figure}

\begin{center}
    {\LARGE \textbf{\textcolor{cyan}{GATE 2009, EEE Question Number 13}}}
\end{center}
\vspace{1em}

\begin{multicols}{2}

\noindent\textbf{Abstract} \\[0.5em]
This project demonstrates the behavior of the NAND gate using a Raspberry Pi Pico2W, push buttons as logic inputs, and an LED as output indicator. It validates the concept that NAND is a universal gate, from which all other logic gates can be built.

\vspace{1em}
\noindent\textbf{1. Components}
\begin{table}[H]
\small
\centering
\begin{tabular}{|p{4.2cm}|c|}
\hline
\textbf{Component} & \textbf{Qty} \\
\hline
Raspberry Pi Pico2W & 1 \\
Push Buttons & 2 \\
LED & 1 \\
220$\Omega$ Resistors & 3 \\
Breadboard & 1 \\
Jumper Wires & 10 \\
Laptop with Thonny IDE & 1 \\
\hline
\end{tabular}
\caption*{Table: Components used}
\end{table}

\vspace{1em}
\noindent\textbf{2. Setup Instructions}
\begin{itemize}
    \item Connect button A to GPIO14 with pull-down configuration.
    \item Connect button B to GPIO15 similarly.
    \item Connect the LED anode to GPIO13 through a 220$\Omega$ resistor.
    \item Common ground is shared between both push buttons and the LED cathode.
    \item Use Thonny IDE to upload the NAND logic implementation.
\end{itemize}

\vspace{1em}
\noindent\textbf{3. NAND Gate Logic}
\[
\text{NAND}(A, B) = \overline{A \cdot B}
\]
\[
\text{Output} = 
\begin{cases}
1 & \text{if } A = 0 \text{ or } B = 0 \\
0 & \text{if } A = 1 \text{ and } B = 1
\end{cases}
\]

\end{multicols}

\vspace{1em}
\noindent\textbf{4. Observation Table}
\[
\begin{array}{|c|c|c|}
\hline
A & B & LED (Output) \\
\hline
0 & 0 & 1 \\
0 & 1 & 1 \\
1 & 0 & 1 \\
1 & 1 & 0 \\
\hline
\end{array}
\]

\vspace{2em}
\noindent\textbf{5. Circuit Image} \\
\begin{center}
\includegraphics[width=0.6\linewidth]{out.jpg} 
\end{center}

\vspace{1em}
\noindent\textbf{6. Applications of NAND Gate}
\begin{itemize}
    \item Used to implement all other basic logic gates (AND, OR, NOT).
    \item Fundamental building block in digital electronics and processor design.
    \item Used in flip-flops, counters, and memory circuits like SRAM.
    \item Employed in Arithmetic Logic Units (ALUs) for logical operations.
\end{itemize}

\vspace{1em}
\noindent\textbf{7. Conclusion}

This hardware implementation successfully verifies the theoretical behavior of a NAND gate using Raspberry Pi Pico2W and basic electronic components. NAND, being a universal gate, is a crucial part of digital design. Understanding and demonstrating it physically enhances our knowledge of logic gate operations and their importance in digital circuit design.

\vspace{1em}
\noindent\textbf{8. References}
\begin{itemize}
    \item GitHub Repository: \url{https://github.com/Rithrishareddy/RITHRISHA_FWC/tree/main/Hardware}
    \item Raspberry Pi Pico Documentation
    \item GATE 2009 EEE Question Paper
\end{itemize}

\end{document}
